\documentclass[11pt]{article}
\usepackage[margin=25mm]{geometry}
\usepackage{enumitem}
\usepackage{xcolor}
\usepackage{hyperref}
\usepackage{url}

\hypersetup{
  hidelinks,
  pdftitle={Response to the A1 Administrative Inquiry for S-28-4-002},
  pdfauthor={Kotaro Furukawa}
}
\setlength{\parindent}{0pt}
\setlength{\parskip}{4pt}
\newcommand{\responseitem}[2]{\section*{#1: #2}}

\begin{document}

\begin{center}
{\Large Response to the A1 Administrative Inquiry}\\[4pt]
Paper No. S-28-4-002\\
Final title: \textit{Camerala: An Edge-Side, Browser-Native Framework for Remote Psychophysiological Experiments}
\end{center}

We thank the Editorial Committee for the A1 decision and for identifying the remaining points requiring correction before publication.
We addressed all six required conditions and four additional comments through factual verification and limited textual revision.
The revised claims remain restricted to the tested N=1 deployment and system-integration implementation.

\responseitem{Required Condition 1}{Resolve inconsistencies in the reported experimental conditions}

\textbf{Comment.}
Please verify and unify the inconsistent descriptions of the block/trial structure, display specifications, and use of an external illuminance meter.

\textbf{Response.}
We verified and unified these experimental conditions throughout the manuscript.
The task used six blocks of ten trials per session, yielding 60 trials per session and 600 trials across ten sessions.
The deployment computer was a mouse Computer G-Tune P517G60ZO21CNHWT with a built-in 15.6-inch, 2560$\times$1440-pixel, 165-Hz display.
No external illuminance measurement was performed.

\textbf{Revision.}
The Methodology now states: ``Each session consisted of six 10-trial blocks, with two blocks for each instructional task context (neutral instruction, negative instruction, and positive instruction), as reflected in the exported block plans.''
The hardware table reports the confirmed computer and display specifications.
The environment section and table now state: ``No external illuminance measurement was performed.''
The inconsistent 3-block $\times$ 20-trial, 13.5-inch/approximately 60-Hz, and lux-meter descriptions do not remain in the revised manuscript.

\textbf{Location.}
Methodology, Target Usage Scenario and Exploratory Deployment (p.~5); Table~2 (p.~6); Deployment Environment and Lighting Conditions and Table~3 (pp.~5--7).

\responseitem{Required Condition 2}{Remove ``with rPPG Sensing'' from the title}

\textbf{Comment.}
The title should not imply validated rPPG waveform extraction or heart-rate estimation when the implementation records facial green-channel intensity traces.

\textbf{Response.}
We removed ``with rPPG Sensing'' from the title.

\textbf{Revision.}
The final title is ``Camerala: An Edge-Side, Browser-Native Framework for Remote Psychophysiological Experiments.''
The manuscript describes the current outputs as facial ROI color traces, facial ROI green-channel values, or face-image-derived records.
It uses rPPG terminology only for related work and for the possible subsequent analysis for which those color traces may be used.

\textbf{Location.}
Title and Abstract (p.~1); Introduction (p.~1); System Architecture (p.~3); Conclusion (p.~9).

\responseitem{Required Condition 3}{Remove the percentage-based criterion from RQ1}

\textbf{Comment.}
The post-hoc percentage criterion should not be used as a criterion for successful system operation.

\textbf{Response.}
We removed the 80\% criterion from RQ1 and from all principal completion claims.

\textbf{Revision.}
RQ1 now reads: ``RQ1: Under the tested configuration, did Camerala complete all planned sessions and trials and export timestamped task-event and face-image-derived records using the same browser clock domain?''
The Results answer this question from completion of all 10 sessions and 600 trials, export of both record streams for every session, and their use of \texttt{performance.now()} in the same client-side browser runtime.
The 100--1500 ms range and 541/600 count remain only as trial-level descriptive preprocessing information and are not used as criteria for procedure completion, timestamp accuracy, or general performance.

\textbf{Location.}
Introduction, Research Questions (pp.~1--2); Methodology, Dependent Measures and Pre-processing (p.~6); Results and Table~4 (pp.~6--8); Conclusion (p.~9).

\responseitem{Required Condition 4}{Replace ``physiological signals'' in RQ2}

\textbf{Comment.}
RQ2 should name the records that were actually exported rather than describing them as validated physiological signals.

\textbf{Response.}
We replaced the physiological-signal wording in RQ2 with terminology that directly identifies the exported records.

\textbf{Revision.}
RQ2 now reads: ``RQ2: What between-environment differences are observable in the task-synchronized face-image-derived and sensor-derived exported records?''
The Results and Discussion describe laboratory/home differences only as differences in exported facial ROI green-channel values, motion estimates, and the frame-to-frame green-channel change indicator.
They are interpreted as mixed deployment-context differences because environmental, camera, task, and participant factors were not independently separated.

\textbf{Location.}
Introduction, Research Questions (pp.~1--2); Results, Between-Environment Exported-Record Differences (p.~7); Discussion (p.~8); Conclusion (p.~9).

\responseitem{Required Condition 5}{Define or remove the camera-derived brightness ranges}

\textbf{Comment.}
The previously reported values 105--114 and 102--111 require a reproducible definition and provenance; otherwise, they should be removed.

\textbf{Response.}
We removed the two setup ranges because their calculation procedure and provenance could not be reconstructed from the saved code and logs.

\textbf{Revision.}
The ranges were deleted from the Methodology, environment table, and all summary passages rather than assigned an unsupported definition.
The manuscript now states only that ``No external illuminance measurement was performed.''
The separately implemented frame-to-frame green-channel change indicator remains defined in the System Architecture from its verified calculation.

\textbf{Location.}
System Architecture, Feature Extraction Logic (p.~4); Methodology, Deployment Environment and Lighting Conditions and Table~3 (pp.~5--7).

\responseitem{Required Condition 6}{Remove the \texttt{mean\_quality = 1.0} result}

\textbf{Comment.}
The exported-window \texttt{mean\_quality} value should not be reported as a quality or tracking-success result because of the structure of the export path.

\textbf{Response.}
We removed the \texttt{mean\_quality} result from the Results, tables, summary passages, RQ criteria, and conclusion.

\textbf{Revision.}
The numeric 299-window/1.0 result is no longer reported.
Only the following structural limitation remains: ``The normal export path stored frame records only when FaceMesh landmarks were returned; therefore, the saved records cannot be used to estimate an overall landmark-detection success rate across all camera frames.''

\textbf{Location.}
Discussion, Limitations (p.~8).

\responseitem{Other Comment 1}{Reorganize the Introduction}

\textbf{Comment.}
The Introduction should begin from the research background and gap rather than a repeated list of claims that the system does not make.

\textbf{Response.}
We reorganized the Introduction into a background-to-scope progression.

\textbf{Revision.}
The four opening paragraphs now address: (1) synchronization and raw-video handling in remote task experiments; (2) existing browser task/rPPG systems and the integration gap; (3) Camerala's integrated browser-side implementation; and (4) the tested seated-task scope.
The detailed limitations are concentrated in the Limitations subsection.

\textbf{Location.}
Introduction (pp.~1--2); Discussion, Limitations (p.~8).

\responseitem{Other Comment 2}{Report camera settings as future reproducibility requirements}

\textbf{Comment.}
Actual negotiated camera resolution, effective frame rate, auto-exposure state, and auto-white-balance state were not recorded and should not be inferred.

\textbf{Response.}
We retained these values as unlogged and added the requested future-work statement.

\textbf{Revision.}
The hardware table distinguishes requested camera constraints from unlogged negotiated values.
The Discussion now states: ``Future evaluations should log the actual negotiated camera resolution and effective frame rate, together with relevant camera-control settings such as auto-exposure and auto-white-balance, to improve reproducibility and interpretation.''

\textbf{Location.}
Table~2 (p.~6); Discussion (p.~8).

\responseitem{Other Comment 3}{Remove revision-process meta-language}

\textbf{Comment.}
Revision-process wording should not appear in the manuscript body.

\textbf{Response.}
We removed the revision-process meta-language from the manuscript.

\textbf{Revision.}
The phrase ``rather than inferred from the revision-time environment'' was replaced with the manuscript-facing statement: ``Values unavailable from the deployment records are explicitly marked as not logged.''
A full-text search confirmed that revision-time and reviewer-facing expressions do not remain in the manuscript body.

\textbf{Location.}
Methodology, Hardware and Software Configuration (p.~5).

\responseitem{Other Comment 4}{Reduce repeated negative statements}

\textbf{Comment.}
Repeated negative claims and limitations should be consolidated without strengthening the manuscript's claims.

\textbf{Response.}
We reduced repeated negative statements and consolidated the detailed limitations.

\textbf{Revision.}
The Abstract gives one concise scope limitation, the end of the Introduction defines the study scope once, and the Limitations subsection contains the detailed constraints concerning N=1, reference sensors, camera settings, landmark-returned export records, and browser timing.
The Conclusion now gives a short limited implementation claim without using RT inclusion or landmark-summary values as evidence of system success.

\textbf{Location.}
Abstract and Introduction (pp.~1--2); Discussion, Limitations (p.~8); Conclusion (p.~9).

\section*{Publication Information}

The final source specifies Vol.~28, No.~4, a received date of April 26, 2026, and a revised date of July 17, 2026 using the commands provided by \texttt{ehis.cls}.
No acceptance date was added.

\end{document}
